\section{Discussion}\label{sec:discussion}

\subsection{What the results establish}

Three findings carry over beyond the particular twin and processor used here. The first is that a physical model of the receiver is a better basis for a monitor than a learned model of attacks. A monitor that knows what an honest link can produce needs no attack data, controls its false-alarm rate through validation on honest links, treats loss as a nuisance rather than an anomaly, and reports which physical constraint an alarm violates. On the benchmark it detected more attacks than every comparison method in every attack group except the long-link regime, where the only method with a higher nominal rate achieved it by alarming on \FALongQBERstatic\,\% of honest long links.

The second finding is that detectability is a property of the attack and the telemetry, and can be computed. The exponent $\Dstar$ separates three situations that a test-set accuracy conflates: an attack that a monitor misses although the information is there, an attack that no monitor can catch within the available time, and an attack that leaves no trace at all. Reporting attack strength in units of $\Dstar$, or of the delay bound $\log(1/\alpha)/(N\Dstar)$, would make results from different systems comparable.

The third finding concerns practice. On real hardware, every monitor tested, including the plain QBER monitors that systems run today, raised alarms on honest links because the reference statistics had been acquired under different operating conditions. Monitoring is only as good as its reference, and the reference must be acquired in operation.

\subsection{Relation to QBER monitoring and to device-independent protocols}

On the qubit hardware, where nearly every disturbance shows up as a change of an error rate, the per-basis QBER test was close to optimal and \qtrace{} matched it. On the optical twin the situation is different: loss, multi-photon pulses, two detectors and three intensities give the receiver statistics a structure that QBER summarises poorly, and the attacks that matter most in practice (PNS, blinding, time-shift, randomness failures) exploit exactly that structure. The advantage of the honest-manifold approach therefore grows with the physical richness of the telemetry. Measurement-device-independent protocols \cite{LoCurtyQi2012} remove detector side channels by construction; for the large installed base of prepare-and-measure systems, a monitor of the kind described here is a low-cost complement that needs no hardware change.

\subsection{Deployment considerations}

The computational cost is modest: calibration takes well under a second per link and one monitoring step requires a one-dimensional Newton iteration and a few small matrix operations per window. The commissioning windows must be attack-free, which is the same trust assumption made when a system is first calibrated. Re-commissioning after maintenance is necessary, and the moving-reference channels make the monitor tolerant to slow drift between commissionings. An alarm need not abort the protocol; it can be used to exclude the affected block from key distillation, as proposed in the operational framework of Mohamed and Al-Kuwari \cite{MohamedAlKuwari2026}. Finite-size key rates \cite{Tomamichel2012,Lim2014} are computed from the same statistics and can be combined with the alarm decision.

\subsection{Limitations}\label{sec:limitations}

The study has limitations that bound its claims. (i) The simulation results depend on the twin and its parameter ranges, which we chose to be realistic but which cannot represent every system; the twin contains mismatch effects but is still a model. (ii) False-alarm control is empirical, through validation quantiles, rather than exact; windows within a link are not independent, and the run-level false-alarm rate depends on the monitoring horizon. (iii) The detectability bound assumes a stationary honest link and is asymptotic; under drift the achievable delay is larger. (iv) The attacks are physically motivated but not adaptive: an eavesdropper who knows the monitor would minimise $\Dstar$, and the space of such attacks has not been explored. (v) The hardware study emulates the channel with superconducting qubits, used a single processor and a single run, has 12 to 16 link-epochs per probe, and its condition-matched analysis was defined after the pre-registered one failed. A follow-up run with timing-matched commissioning circuits is needed to confirm the matched results prospectively. (vi) The afterpulse probability is taken from a specification; in systems where it is unknown it must be calibrated jointly with the background, which is poorly identified without loss variation during commissioning.

\section{Conclusion}\label{sec:conclusion}

We have shown that implementation attacks on decoy-state QKD can be detected without attack data by testing receiver telemetry against the manifold of statistics that an honest, lossy link can produce. The construction combines an exact count-level receiver model, an efficient treatment of the transmittance nuisance, physically structured score tests and drift-robust references. It comes with a detector-independent exponent that bounds what any telemetry monitor can achieve and that identifies attacks no monitor can see. On a benchmark with deliberate model mismatch the monitor detected substantially more attacks than industry monitors, unsupervised detectors and a supervised classifier at the same false-alarm rate, with delays close to the information-theoretic bound. A hardware study with physically executed eavesdroppers confirmed the probe physics, demonstrated low-rate attack detection with a condition-matched reference, and showed that a reference acquired under different operating conditions defeats every monitor, including the QBER monitors in use today. Natural next steps are an extension of the bound to drifting links, the study of adaptive attacks that minimise $\Dstar$, finite-key integration, and validation on an optical testbed.
